% META: {"id": "1NSI-PM-AM-EXTRAIT", "chapitre": "1NSI-PROJET-METHODES", "type_objet": "amenagee", "capacites": ["BO-PREAMBULE-DEMARCHE-DE-PROJET", "BO-PREAMBULE-COMPETENCES-METHODE"], "capacites_codes": ["C2", "C3"], "derive_de": ["1NSI-PM-EX-002", "1NSI-PM-EX-003"], "profil_amenagement": ["SEQUENCAGE", "ALLEGEMENT_ENONCE", "REPONSE_FERMEE", "AERATION", "AMORCE_FOURNIE", "RETRAIT_DOUBLE_TACHE"], "mode_creation": "adaptation", "sources_inspiration": [], "status": "generated"}
\section*{Version aménagée --- Extrait Projet et méthodes}
\textit{Consignes séquencées, mise en page aérée, énoncés allégés.}

\bigskip

\subsection*{Exercice 1 --- Reconnaître un bon jalon (C2)}

Une équipe de trois élèves développe un petit jeu. Voici quatre propositions.

\bigskip

\textbf{Étape 1.} Pour chaque proposition, coche si c'est un vrai jalon.

\medskip
\textit{Rappel :} un jalon a une \textbf{date} et un \textbf{état vérifiable}.

\bigskip

\begin{center}
\begin{tabular}{|l|c|}
\hline
\textbf{Proposition} & \textbf{Vrai jalon ?} \\
\hline
« Avancer sur les graphismes » & \dotfill \\
\hline
« Le 12 mars, le joueur peut se déplacer » & \dotfill \\
\hline
« Travailler dur cette semaine » & \dotfill \\
\hline
« Le 19 mars, les tests du score passent » & \dotfill \\
\hline
\end{tabular}
\end{center}

\bigskip

\textbf{Étape 2.} Pourquoi « Avancer sur les graphismes » n'est-il pas un jalon ?

\medskip
\begin{center}
\fbox{on ne peut pas dire s'il est atteint} \qquad \fbox{c'est trop difficile}
\end{center}

\bigskip

\textbf{Étape 3.} Le rendu est le $30$ mars. Le dernier jalon doit-il tomber ce jour-là ?

\medskip
\begin{center}
\fbox{Oui} \qquad \fbox{Non, il faut garder une marge}
\end{center}

\bigskip
\bigskip

\subsection*{Exercice 2 --- Lire une erreur (C3)}

Ce programme plante :

\begin{python}
notes = {"Alice": 15, "Bob": 12}

def afficher_note(notes, eleve):
    print(notes[eleve])

afficher_note(notes, "Chloe")
\end{python}

\medskip

\begin{console}
Traceback (most recent call last):
  File "programme.py", line 6, in <module>
    afficher_note(notes, "Chloe")
  File "programme.py", line 4, in afficher_note
    print(notes[eleve])
KeyError: 'Chloe'
\end{console}

\bigskip

\textbf{Étape 1.} Complète le tableau. Lis le message de bas en haut.

\bigskip

\begin{center}
\begin{tabular}{|l|c|}
\hline
\textbf{Question} & \textbf{Réponse} \\
\hline
Type de l'erreur & \dotfill \\
\hline
Ligne où elle se produit & \dotfill \\
\hline
Clé qui pose problème & \dotfill \\
\hline
\end{tabular}
\end{center}

\bigskip

\textbf{Étape 2.} Quelle est la cause ? Entoure.

\medskip
\begin{center}
\fbox{« Chloe » n'est pas dans le dictionnaire} \qquad \fbox{le dictionnaire est vide}
\end{center}

\bigskip

\textbf{Étape 3.} Complète la correction. Une seule ligne à ajouter.

\begin{python}
def afficher_note(notes, eleve):
    if eleve ______ notes:
        print(notes[eleve])
    else:
        print("eleve inconnu")
\end{python}

\medskip
Le mot manquant : \fbox{\texttt{in}} \qquad ou \qquad \fbox{\texttt{==}}
